\documentclass[11pt]{article}
\usepackage[margin=0.7in]{geometry}
\usepackage{amsmath,amssymb}
\usepackage{enumitem}
\usepackage{tikz}
\usepackage{pgfplots}
\pgfplotsset{compat=1.18}
\usepackage{fancyhdr}
\usepackage{parskip}
\setlist[enumerate]{leftmargin=*, itemsep=6pt}
\pagestyle{fancy}
\fancyhf{}
\fancyfoot[L]{\small\textit{en}Vision\textsuperscript{\textregistered} Algebra 2}
\fancyfoot[C]{\small 4-3 Lesson Quiz, short form}
\fancyfoot[R]{\small\thepage}
\renewcommand{\headrulewidth}{0pt}
\newcommand{\choice}[1]{\tikz[baseline=-0.6ex]\node[draw, circle, inner sep=0.6pt, minimum size=3.3mm, font=\scriptsize\bfseries] {#1};\;}
\newcommand{\checkbox}{$\square$\;}
\pgfplotsset{quizgrid/.style={
  axis lines=middle, grid=both, grid style={gray!45},
  ticklabel style={font=\scriptsize, fill=white, inner sep=1pt},
  xlabel={$x$}, ylabel={$y$}, axis line style={thick}, clip=true}}
\renewcommand{\arraystretch}{1.6}

\begin{document}

\noindent\textbf{Name}\; \rule{6cm}{0.4pt} \hfill Period \rule{1.2cm}{0.4pt} \quad Date \rule{2.2cm}{0.4pt}

\vspace{4pt}\noindent\rule{\textwidth}{0.6pt}

\vspace{2pt}\noindent{\Large\textbf{4-3 \; Lesson Quiz: Multiplying and Dividing Rational Expressions}} \hfill {\small short form}

\vspace{4pt}\noindent\rule{\textwidth}{0.6pt}

{\small Question numbers match the Savvas Lesson Quiz. Questions 1 and 3 are not on this form. Show your factoring for every question.}

\vspace{4pt}
\begin{minipage}[t]{0.52\textwidth}
\textbf{2.} Simplify $\dfrac{(x+5)(x^2-25)}{(x+5)^2(x-5)^2}$. What is the domain?

\vspace{6pt}
\begin{itemize}[leftmargin=1.2em, itemsep=6pt]
\item[] \choice{A} $\dfrac{x+5}{x-5}$; all real numbers except $-5$ and $5$
\item[] \choice{B} $\dfrac{1}{x-5}$; all real numbers except $-5$ and $5$
\item[] \choice{C} $\dfrac{1}{x+5}$; all real numbers except $0$
\item[] \choice{D} $\dfrac{x-5}{x+5}$; all real numbers except $-5$
\end{itemize}
\end{minipage}\hfill
\begin{minipage}[t]{0.44\textwidth}
\textbf{Work:} factor, then name the values that make the original denominator zero.

\vspace{4pt}\rule{\linewidth}{0.4pt}

\vspace{14pt}\rule{\linewidth}{0.4pt}

\vspace{14pt}\rule{\linewidth}{0.4pt}

\vspace{14pt}\rule{\linewidth}{0.4pt}
\end{minipage}

\vspace{10pt}\noindent\rule{\textwidth}{0.4pt}

\vspace{6pt}
\begin{minipage}[t]{0.52\textwidth}
\textbf{4.} What is the quotient of $\dfrac{9-x^2}{3x}$ and $\dfrac{x^2+6x+9}{3x}$?

\vspace{6pt}
\begin{tabular}{@{\hspace{1.2em}}l@{\hspace{2.5em}}l}
\choice{A} $\dfrac{x-3}{3-x}$ & \choice{B} $\dfrac{x-3}{x+3}$ \\[10pt]
\choice{C} $\dfrac{3-x}{x+3}$ & \choice{D} $\dfrac{x+3}{x-3}$
\end{tabular}
\end{minipage}\hfill
\begin{minipage}[t]{0.44\textwidth}
\textbf{Work:} write the division as a multiplication by the reciprocal, then factor.

\vspace{4pt}\rule{\linewidth}{0.4pt}

\vspace{14pt}\rule{\linewidth}{0.4pt}

\vspace{14pt}\rule{\linewidth}{0.4pt}

\vspace{14pt}\rule{\linewidth}{0.4pt}
\end{minipage}

\vspace{10pt}\noindent\rule{\textwidth}{0.4pt}

\vspace{6pt}
\begin{minipage}[t]{0.52\textwidth}
\textbf{5.} The area of a rectangle is $\dfrac{x^2-4}{2x}$ in.$^2$ and its length is $\dfrac{(x+2)^2}{2}$ in. What is the width in inches?

\vspace{6pt}
\begin{tabular}{@{\hspace{1.2em}}l@{\hspace{2.5em}}l}
\choice{A} $\dfrac{x(x+2)}{x-2}$ & \choice{B} $\dfrac{x+2}{x(x+2)}$ \\[10pt]
\choice{C} $\dfrac{x-2}{x(x-2)}$ & \choice{D} $\dfrac{x-2}{x(x+2)}$
\end{tabular}
\end{minipage}\hfill
\begin{minipage}[t]{0.44\textwidth}
\textbf{Work:} width $=$ area $\div$ length.

\vspace{4pt}\rule{\linewidth}{0.4pt}

\vspace{14pt}\rule{\linewidth}{0.4pt}

\vspace{14pt}\rule{\linewidth}{0.4pt}

\vspace{14pt}\rule{\linewidth}{0.4pt}
\end{minipage}

\vfill
\noindent\textbf{Scoring:} Question 2 and Question 4 are 10 points each. Question 5 is a bonus of up to 5 points. Total out of 20.

\newpage
\noindent{\LARGE\textbf{4-3 Lesson Quiz, short form: answer key}}

\vspace{4pt}\noindent\rule{\textwidth}{0.6pt}

\begin{enumerate}
\item[\textbf{2.}] \textbf{B, $\dfrac{1}{x-5}$; all real numbers except $-5$ and $5$.} $x^2-25=(x+5)(x-5)$, so the expression is $\dfrac{(x+5)(x+5)(x-5)}{(x+5)^2(x-5)^2}=\dfrac{1}{x-5}$. The domain comes from the original denominator $(x+5)^2(x-5)^2$: $x\neq-5$ and $x\neq5$. Choices A, C and D are not equivalent to the expression; C and D also state the wrong domain (C leaves out only $0$, D leaves $5$ in). Savvas quiz item 2, DOK 2, skill A85; taught by Example 2. \textbf{10 points:} 6 for B, 4 for the work (both factorizations shown and the excluded values named from the original denominator).
\item[\textbf{4.}] \textbf{C, $\dfrac{3-x}{x+3}$.} $\dfrac{9-x^2}{3x}\cdot\dfrac{3x}{x^2+6x+9}=\dfrac{(3-x)(3+x)\cdot 3x}{3x\,(x+3)(x+3)}=\dfrac{3-x}{x+3}$ (domain, not asked: $x\neq0,\,-3$). Choice B is the sign slip $x-3$ for $3-x$ (they are opposites); D is the reciprocal of B. Savvas quiz item 4, DOK 2, skill A85; taught by Example 5. \textbf{10 points:} 6 for C, 4 for the work (the reciprocal written, $9-x^2$ and $x^2+6x+9$ factored).
\item[\textbf{5.}] \textbf{D, $\dfrac{x-2}{x(x+2)}$.} Width $=$ area $\div$ length $=\dfrac{x^2-4}{2x}\cdot\dfrac{2}{(x+2)^2}=\dfrac{(x+2)(x-2)\cdot2}{2x\,(x+2)(x+2)}=\dfrac{x-2}{x(x+2)}$. Savvas quiz item 5, DOK 2, skill A85; taught by Example 5 (division), rehearsed by the packet's Practice (TE Practice 37, base of a parallelogram $=$ area $\div$ height). \textbf{Bonus 5:} 5 for D with the division set up as area times the reciprocal of the length; 2 for D with no work.
\end{enumerate}

\textbf{Items left off this form} (available for a retake or Form B): Savvas item 1, simplify $\dfrac{x^2-2x}{(x-2)(x+3)}$ (answer C, $\dfrac{x}{x+3}$; DOK 2, A85; Example 2); Savvas item 3, simplify $\dfrac{ab^2-a^2b}{ab}\cdot\dfrac{a^2b^2}{a-b}$ (answer A, $-a^2b^2$; $a\neq0$, $b\neq0$; DOK 2, A85; Example 3, cut in the fast packet). The TE's choice A omits $a\neq b$, which the original denominator $a-b$ requires; accept A and mention it.

\textbf{IXL replacement:} SmartScore 70 or above on 37N, 79Q and MG2 replaces this quiz with 20/20.

\end{document}
